%=====================================================================
%  ch01 —— 原文 297–496：题名、摘要、§1 引言、致谢
%=====================================================================

\section{引言}

\subsection{} 群的范畴作用理论与局部几何朗兰兹对应，是现代表示论中两个重要的角色。
过去二十年里，它们发展迅速：Beilinson、Drinfeld、Frenkel、Gaitsgory 作出了决定性的贡献，
更晚近则有 Beraldo、Raskin、Yang。

无论从历史看，还是从主要应用看，范畴作用与局部几何朗兰兹都同李代数的表示论深深交织。
这两个联系，本书会反复强调。

一方面，范畴作用的理论为处理约化李代数的表示论提供了高效的框架。这也是进入范畴作用理论的
一条便利路径，本书正是由此入手。

另一方面，也是更重要的方面，范畴作用的形式体系与局部几何朗兰兹的视角，对仿射李代数及相关
顶点代数的表示论发展不可或缺。它们既提供了解决旧猜想的工具，也催生了新的猜想。这些进展将在
本书后半部分讲到。

\subsection{} 这一圈想法有一个美妙之处：尽管表面上用到了现代代数与几何中相当精致的工具，
只要讲法得当，非专家也能迅速吸收其中的基本结果、技巧与视角。我们希望本书能为非专家读者
提供这样一个人口。

为此，我们给自己定下的目标是：以可读的方式综述这一领域，并补足文献中已有的材料。出于这一
考虑，我们选择着重讲类比、动机与历史前身，而不去对所涉材料作完全严格的展开。

\subsection{} 追求这样一个目标，立刻会遇到两个问题。

第一个问题是：假定读者已经知道多少，又该把多少内容讲得细一些。我们假定读者对单李群与单李
代数的表示论、代数几何、同调代数的基础有所了解，但不超过这些科目入门课程的内容。当然，
这些预备也可以换成一定的数学成熟度，或者至少是愿意把若干断言先信下来。特别地，我们不假定
读者熟悉几何表示论的工具，例如 D-模与反常层、高阶代数，以及导出与无穷维代数几何。需要什么，
我们就在用到时发展什么。这样做的同时，我们仍力求以最少的枝节传达基本想法，希望读者借此
获得方向感，再去攻读有关文献。

第二个问题是：强调哪些想法、方法与直觉，必然受篇幅限制，也受综述者本人偏好与局限的限制。
因此，对于许多贡献被略去或强调不足的数学家与物理学家，我们在此致歉。先说一点：本书与三维、
四维超对称规范理论之间极其丰富的联系，只在末尾略作提及。

\subsection{} 我们努力使全书的材料复杂度单调递增。读者可以从头读起，也可以跳到觉得陌生
或有趣之处，从那里往下读。我们希望前几节，特别是第 \ref{s:bbloc} 节，能作为非专家进入
几何表示论的一篇轻量导引。至于最后几节，特别是第 \ref{s:affbbloc} 节，我们也预计，要领会
其中每一个细节，需要一定程度的预备知识。

本书的具体安排如下。第 \ref{s:repthryfingroup} 节回顾有限群表示论的若干基本方面，为后面
遇到的范畴对应物作铺垫。第 \ref{s:bbloc} 节回顾 Beilinson--Bernstein 局部化定理，并说明
由此得到的等变 D-模与该李代数的 Harish-Chandra 模之间的对应，如何引出范畴群作用的想法。
第 \ref{s:catrepgrp} 节讨论有限维代数群的范畴群作用理论的基础，并花一些篇幅解释现代同调代数
中若干必要的工具与构造，好让读者能够上手。要走到局部几何朗兰兹，必须从代数群的情形过渡到
环路群。因此，第 \ref{s:catloop} 节建立环路群的范畴表示论的基础，也就是检讨进入无穷维情形
后哪些地方要变。最后，第 \ref{s:affbbloc} 节讲局部几何朗兰兹对应，较详细地讨论其中已知与
猜想的部分，并着重介绍它在临界水平仿射李代数的局部化理论中的一个重要应用。附录
\ref{s:frmfun} 还为非专家提供了一份简明导引，讲 D-模、Riemann--Hilbert 对应与函数-层对应。

\subsection{} 除书中各处给出的文献指引外，读者还可参阅关于局部几何朗兰兹的两部出色综述
\cite{frenkellocalnotes} 与专著 \cite{frenkelloop}，以及讲义 \cite{paris-notes}；关于整体
图景，\cite{frenkelglobalnotes} 是一份通俗的导引。

\vspace{.5cm}

\noindent{\bf 致谢。} 感谢 Pramod Achar、Dima Arinkin、Dario Beraldo、Roman Bezrukavnikov、
Alexander Braverman、Justin Campbell、Harrison Chen、Thomas Creutzig、Davide Gaiotto、
Tony Feng、Edward Frenkel、Igor Frenkel、Dennis Gaitsgory、Ivan Losev、Sam Raskin、
Daniel Sage、Germ\'an Stefanich、Jonathan Wang、David Yang、Zhiwei Yun、Xinwen Zhu，
与他们的讨论加深了我们对本课题的理解。我们还要深深感谢 Cora Bernard、Apoorva Khare
以及匿名审稿人：他们仔细阅读了本文早先的稿本，并提出了有益的意见。

作者的工作由美国国家科学基金会（NSF）的博士后资助支持，资助号 2103387。
